\documentclass[11pt]{article}
\usepackage{fontspec}
\setmainfont{Times New Roman}
\setsansfont{Arial}
\setmonofont{Consolas}
\usepackage{amsmath,amssymb,mathtools}
\usepackage{geometry}
\usepackage{microtype}
\geometry{a4paper,margin=2.35cm}
\setlength{\parindent}{1.25em}
\setlength{\parskip}{0pt}
\makeatletter
\renewcommand\footnotesize{\@setfontsize\footnotesize{7.7}{8.7}\selectfont}
\makeatother
\emergencystretch=100em
\allowdisplaybreaks
\sloppy
\hbadness=10000
\vbadness=10000
\hfuzz=2pt
\vfuzz=10pt
\raggedbottom

\begin{document}

% Unit: Paper31 Section04 entry 6 v001. Direct Interslavic Latin translation from the German final-audited source slice, source lines 283--287 / segment P31-S0061, expanded against printed scan/OCR witness page 94. Footnote counter continues after Paper31 Section04 entry 5.
\setcounter{footnote}{26}

\textbf{3. \emph{Slěd elementa vzhodno k klasu.}} Naj $C$ bude matrica, prirěđena elementu $\gamma$ iz $\mathfrak R$ v $R_{\mathfrak a}$. Prěhod od $R_{\mathfrak a}$ k ekvivalentnomu kolcu pokazuje, že determinant $|C|$, i obče $|tE-C|$, gdje $t$ označa neopreděljenu, jest invariant klasa predstavjenj. Tym samym, po punktu 2, to jest jednovrěmenno invariant klasa idealov, to jest klasovy invariant, kako budemo kratko govoriti. Zato pojedine koeficienty po $t$ v $|tE-C|$ sut klasove invarianty; osoblivo koeficient pri $(-t^{s-1})$ budemo nazyvati slědom $\gamma$ vzhodno k klasu; v znakih:
\[
        S_{(\mathfrak a)}(\gamma)=c_{11}+c_{22}+\cdots+c_{ss}.
\]
Koeficient svobodny od $t$, $\pm |C|$, budemo nazyvati normoju $\gamma$ vzhodno k klasu.

\end{document}
